% !TEX TS-program = lualatex+se+debug
% !TEX encoding = UTF-8 Unicode

\documentclass[14pt, twoside, a4paper]{extreport}
% -----------------------------------------------------------------------------------
% ------------------------- USTAWIENIA STRONY--------------------------
% -----------------------------------------------------------------------------------
\usepackage{geometry}
\geometry{
	left=10mm,
	right=10mm,
	top=15mm,
	bottom=15mm,
	marginparwidth=0mm,
	marginparsep=0mm
	}
\newcommand{\wstawplik}[1]{\input{../include/#1}}%skraca kod
\wstawplik{headerofficium}
% elementy zmienne co do headerofficium:
\graphicspath{{../img/}} % katalog obrazów
\grechangestaffsize{20}%rozmiar wyniku GregorioTEX
% koniec 

% debugowanie:
% %\iffalse %*** %\fi %***
% \tracingall przed problematycznym kodem - bardzo dużo informacji
% lualatex+se+debug - dodatkowe opcje: -file-line-error i -halt-on-error 
% \errorcontextlines

\begin{document}
% -----------------------------------------------------------------------------------
% ------------------------- STRONA TYTUŁOWA  --------------------------
\thispagestyle{empty}
\begin{center}
	\includegraphics[width=1.0\linewidth]{singers_okładka.jpg}
	\vfill
	\tytulfont{COMPLETORIUM}
	\vfill
	\podtytulfont{Kompleta niedzielna\\wraz\\z~modlitwami wieczornymi}
	\vfill
	\large\textit{Liber usualis, Solesmes, 1961, p.262}
\end{center}%
\newpage
%\thispagestyle{empty}
\wstawplik{aperi_domine}
\setcounter{page}{1}
\fancyhead{}% wyczyść wszystkie pola nagłówka
\pagestyle{plain}% nagłówek i stopka pusta - UWAGA: przedefiniowany styl przez polecenie fancypagestyle w header.tex	
%\setlength{\parindent}{0pt}
\gresetbracerendering{font}
\sectionspace
\separator
\sekcja{Incipit}{Wszyscy wstają. Lektor wychodzi na środek, przyklęka i skłaniając się w kierunku ołtarza prosi o błogosławieństwo. W~obecności kapłana: ``Jube domne benedicere''~w jego kierunku. Błogosławieństwo - odpowiada kapłan lub pod jego nieobecność hebdomadariusz.}
\gresetinitiallines{0}\gabcsnippet{(c3)<c><sp>V/</sp>.</c> Ju([alt:Lector]h)be(h) Do(h)mi(h)ne(g'_) be(h)ne(h)di(h)ce(d)re(f.::)}
\vspace{-5mm}
\tlumaczR{\Vbar. Pobłogosław, ojcze.}\segmentspace% 
\gresetinitiallines{0}\gabcsnippet{(c3)Noc([alt:Hebdomadarius]h)tem(h) qui(h)ét(h)am(h) et(h) fi(h)nem(g) per(f)féc(h)tum(h.;) con(h)cé(h)dat(h) no(h)bis(h) Do(h)mi(h)nus(h) o(h)mni(h)po(d)tens(d.::) <c><sp>R/</sp>.</c> A([alt:Chorus]g.)men(h.::)}
\tlumaczR{Noc spokojną i koniec doskonały niech nam da Pan wszechmogący. \Rbar. Amen.}\segmentspace% 
\gresetinitiallines{1}\grecommentary{1 Petr. 5, 8-9}\gabcsnippet{(c3) Frá([alt:Lector]h)tres:(h.)(,) Só(h)bri(h)i(h) e(h)stó(h)te(h), et(h) vi(h)gi(h)lá(h)te:(f.;) qui(h)a(h) a(h)dver(h)sá(h)ri(h)us(h) vés(h)ter(h) diá(h)bo(h)lus,(h'_)(,) tam(h)quam(h) lé(h)o(h) rú(h)gi(h)ens,(h'_)(,) cír(h)cu(h)it(h), quaé(h)rens(g) quem(f) dé(h)vo(h)ret:(h.;) cú(h)i(h) re(h)sí(h)sti(h)te(h) fór(h)tes(h) in(h) fí(h.)de.(d.:) Tu(h) au(g)tem(f) Dó(h)mi(h)ne(h._), mi(h)se(h)ré(h)re(h) nó(h.)bis.(d.::) <c><sp>R/</sp>.</c> De([alt:Chorus]h)o(h) grá(h)ti(d)as.(d.::)}
\tlumaczR{Bracia: Trzeźwymi bądźcie, a czujcie: boć przeciwnik wasz djabeł jako lew ryczący krąży, szukając, kogoby pożarł: któremu się sprzeciwiajcie mocni w wierze. \Vbar. Ty zaś, Panie, zmiłuj się nad nami. \Rbar. Bogu niech będą dzięki.}\segmentspace% 
\gresetinitiallines{0}\gabcsnippet{(c3) <c><sp>V/</sp>.</c> Ad(h[alt:Hebdomadarius])ju(h)to(h)ri(h)um(h) nos(h)trum(h) <c><v>\grecross</v></c> in(h) no(h)mi(h)ne(h) Do(h)mi(f)ni(f) (::) <c><sp>R/</sp>.</c> Qui fe(h[alt:Chorus])cit(h) cae(h)lum(h) et(h) ter(h)ram.(f') (::)}
\tlumaczR{\Vbar. Wspomożenie nasze \grecross~w imieniu Pana. \Rbar. Który stworzył niebo i ziemię.}
\newpage
\rubryka{Odmawia się \emph{Pater noster} po cichu, oraz czyni rachunek sumienia. W przypadku obecności kapłana, może on poprowadzić rachunek sumienia, wówczas pierwszy odmawia  \emph{Confiteor}, a~następnie wszyscy (analogicznie jak przy modlitwach u stopni ołtarza).}
\BeginParallel
\latN{\inicjal{P}{ater} noster, qui es in cælis, sanctificétur nomen tuum: advéniat regnum tuum: fiat volúntas tua, sicut in cælo et in terra. Panem nostrum quotidiánum da nobis hódie: et dimítte nobis débita nostra, sicut et nos dimíttimus debitóribus nostris: et ne nos indúcas in tentatiónem: sed líbera nos a malo. Amen.}\pol{\inicjal{O}{jcze} nasz, któryś jest w niebie, święć się imię Twoje: przyjdź królestwo Twoje: bądź wola Twoja, jako w niebie tak i na ziemi. Chleba naszego powszedniego daj nam dzisiaj: i odpuść nam nasze winy, jako i my odpuszczamy naszym winowajcom: i nie wódź nas na pokuszenie: ale nas zbaw ode złego. Amen.}%

\latN{\inicjal{C}{onfíteor} Deo omnipoténti, beátæ Maríæ semper Vírgini, beáto Michaéli Archángelo, beáto Joánni Baptístæ, sanctis Apóstolis Petro et Paulo, et ómnibus Sanctis, quia peccávi nimis, cogitatióne, verbo et ópere: mea culpa, mea culpa, mea máxima culpa. Ídeo precor beátam Maríam semper Vírginem, beátum Michaélem Archángelum, beátum Joánnem Baptístam, sanctos Apóstolos Petrum et Paulum, et omnes Sanctos, oráre pro me ad Dóminum Deum nostrum.}\pol{\inicjal{S}{powiadam} się Bogu wszechmogącemu, Najświętszej Maryi zawsze Dziewicy, świętemu Michałowi Archaniołowi, świętemu Janowi Chrzcicielowi, świętym Apostołom Piotrowi i Pawłowi, i wszystkim Świętym, że bardzo zgrzeszyłem, myślą, mową i uczynkiem: bijąc się w piersi moja wina, moja wina, moja bardzo wielka wina. Przeto błagam Najświętszą Maryję zawsze Dziewicę, świętego Michała Archanioła, świętego Jana Chrzciciela, świętych Apostołów Piotra i Pawła, i wszystkich świętych, o modlitwę za mnie do Pana Boga naszego.}%

\latN{\inicjal{M}{isereátur} nostri omnípotens Deus, et dimíssis peccátis nostris, perdúcat nos ad vitam ætérnam. Amen.}\pol{\inicjal{N}{iech} się zmiłuje nad nami Bóg wszechmogący, a odpuściwszy nam grzechy, doprowadzi nas do życia wiecznego. Amen.}%

\latN{\inicjal{I}{ndulgéntiam}, \grecross~absolutiónem et remissiónem peccatórum nostrórum tríbuat nobis omnípotens et miséricors Dóminus. Amen.}\pol{\inicjal{P}{rzebaczenia}, \grecross~rozgrzeszenia i odpuszczenia grzechów naszych, niech nam udzieli Pan wszechmogący i miłosierny. Amen.}%
\EndParallel
\rubryka{Czyni się znak krzyża na klatce piersiowej.}
\gresetinitiallines{0}\gabcsnippet{(c3) <c><sp>V/</sp>.</c> Con(h[alt:Hebdomadarius])vér(h)te(h) nos(h) <c><v>\grecross</v></c> Déus(h) sa(h)lu(h)tá(g)ris(f) nó(h)ster(h) (::) <c><sp>R/</sp>.</c> Et(h[alt:Chorus]) a(h)vér(h)te(h) í(h)ram(h)  tú(h)am(g) a(f) nó(h)bis.(h) (::)}
\tlumaczR{\Vbar. Nawróć nas \grecross~Boże, nasz zbawco. \Rbar. I odwróć swój gniew od nas.}\segmentspace% 
\gresetinitiallines{0}\gabcsnippet{(c3) <c><sp>V/</sp>.</c> De(h)us(h'_) <c><v>\grecross</v></c> (,) in(h) ad(h)ju(h)tó(i)ri(h)um(h) me(h)um(h'_) in(h)tén(g)de.(h.) (::) <c><sp>R/</sp>.</c> Dó(h)mi(h)ne(h'_) (,) ad(h) ad(h)ju(h)ván(h)dum(h) me(h'_) fe(h)stí(g)na.(h.) (:) Gló(h)ri(h)a(h) Pa(h)tri,(h) et(h) Fí(h)li(h)o,(h'_) (,) et(h) Spi(h)rí(h)tu(h)i(h) San(g)cto.(h.) (:) Sic(h)ut(h) e(h)rat(h) in(h) prin(h)cí(h)pi(h)o,(h) et(h) nunc,(h) et(h) sem(h)per,(h.) (,) et(h) in(h) saé(h)cu(h)la(h) sae(h)cu(h)ló(h)rum.(h) A(g)men.(h.) (;) Al(h)le(i)lú(hg~){ia}.(g.) (::) <i>Vel :</i>() Laus(h) ti(h)bi(h) Dó(h)mi(h)ne(h'_) Rex(h) ae(h)tér(h)nae(i) gló(h')ri(h)ae.(g.) (::)}
\tlumaczR{\Vbar. Boże, \grecross~wejrzyj ku wspomożeniu memu. \Rbar. Panie, pośpiesz ku ratunkowi memu. Chwała Ojcu i Synowi, i Duchowi Świętemu. \newline Jak była na początku, teraz i zawsze, i na wieki wieków. Amen. Alleluja.~\emph{Lub:} Chwała Tobie Panie, Królu wiecznej chwały.}\newpage
\rubryka{Po zaintonowaniu antyfony, oficjant rozpoczyna pierwszy wers psalmu, wszyscy siadają, po czym drugi chór wspólnie śpiewa kolejny wers. Drugi psalm intonuje kantor (wstając), trzeci ponownie oficjant.}
\begin{center}
\begin{tabularx}{0.8\textwidth}{ p{150pt} p{50pt} p{150pt} p{50pt}}
 \small\emph{Per Annum} &  & \small\emph{Tempore Paschali} \\ 
\begin{minipage}{150pt}\gresetinitiallines{1}\greannotation{Ant.}\greannotation{8. G}\gabcsnippet{(c4) MI(g)se-(g')ré-(h)re(gf) (::)}\end{minipage}  & & \begin{minipage}{150pt}\gresetinitiallines{1}\greannotation{Ant.}\greannotation{8. G}\gabcsnippet{(c4) AL(g)le-(g')lú-(h)ia(gf) (::)}\end{minipage}  & \\
\tlumaczR{Zmiłuj się} & & &
\end{tabularx}
\end{center}
\vspace{-12mm}
\subsekcja{Psalm 4}
\vspace{-2mm}
\gresetinitiallines{0}\gabcsnippet{(c4)<c><v>\footnotesize{1. }</v></c>Cum(g) in(h)vo(j)cá(j)rem(j) ex(j)au(j)dí(j)vit(j) me(j) De(j)us(j) jus(j)tí(j)ti(j)æ(j) <b>me</b>(k jr)æ:(j.) <c><v>*</v></c> (:) in(j) tri(j)bu(j)la(j)ti(j)ó(j)ne(j) di(j)la(j)<i><ul>tás</i></ul>(i)<i><ul>ti</i></ul>(j) <b>mi</b>(h gr)hi.(g.) (::)
<i>Flex:</i>()  pec(j)cá(j)re: <c><v>†</v></c>(h. j j j  ::)}
\vspace{-5mm}
\tlumaczR{Kiedym wzywał, wysłuchał mię Bóg sprawiedliwości mojéj:~\red{*} w uciśnieniu rozprzestrzeniłeś mi.}\segmentspace
\columnratio{0.7}
\setlength{\columnseprule}{0.4pt}
\setlength{\columnsep}{.02\textwidth}
\BeginParallel
\latN{\mred{2.}~Miserére \textbf{me}i,~\red{*}  et exáudi orati\underline{\textit{ó}}\underline{\textit{nem}} \textbf{me}am.}\polPs{Zmiłuj się nademną,~\red{*} a wysłuchaj modlitwę moję.}%

\latN{\mred{3.}~Fílii hóminum, úsquequo gravi \textbf{cor}de?~\red{*}  ut quid dilígitis vanitátem et quǽri\underline{\textit{tis}} \underline{\textit{men}}\textbf{dá}cium?}\polPs{Synowie człowieczy, pókiż ciężkiego serca?~\red{*}  przecz miłujecie marność, i szukacie kłamstwa?}%

\latN{\mred{4.}~Et scitóte quóniam mirificávit Dóminus sanctum \textbf{su}um:~\red{*}  Dóminus exáudiet me cum clamáve\underline{\textit{ro}} \underline{\textit{ad}} \textbf{e}um.}\polPs{Wiedzcież, iż dziwnym uczynił Pan świętego swego:~\red{*}  wysłucha mię Pan, gdy zawołam do niego.}%

\latN{\mred{5.}~Irascímini, et nolíte peccáre:~\red{†} quæ dícitis in córdibus \textbf{ves}tris,~\red{*}  in cubílibus vestris \underline{\textit{com}}\underline{\textit{pun}}\textbf{gí}mini.}\polPs{Gniewajcie się, a nie grzeszcie:~\red{†} co mówicie w sercach waszych,~\red{*}  na łóżkach waszych żałujcie.}%

\latN{\mred{6.}~Sacrificáte sacrifícium justítiæ,~\red{†} et speráte in \textbf{Dó}mino.~\red{*}  Multi dicunt: quis osténdit \underline{\textit{no}}\underline{\textit{bis}} \textbf{bo}na?}\polPs{Ofiarujcie ofiarę sprawiedliwości,~\red{†} a ufajcie w Panu.~\red{*}  Wiele ich mówią: Któż nam okazał dobra?}%

\latN{\mred{7.}~Signátum est super nos lumen vultus tui, \textbf{Dó}mine:~\red{*}  dedísti lætítiam in \underline{\textit{cor}}\underline{\textit{de}} \textbf{me}o.}\polPs{Naznamionowana jest nad nami światłość oblicza twego, Panie:~\red{*}  dałeś wesele w sercu mojem.}%

\latN{\mred{8.}~A fructu fruménti, vini et ólei \textbf{su}i~\red{*}  mul\underline{\textit{ti}}\underline{\textit{pli}}\textbf{cá}ti sunt.}\polPs{Od urodzaju zboża, wina i oliwy swéj~\red{*}  rozmnożeni są.}%

\latN{\mred{9.}~In pace in id\textbf{íp}sum~\red{*}  dórmiam et \underline{\textit{re}}\underline{\textit{qui}}\textbf{és}cam;}\polPs{W pokoju pospołu~\red{*}  będę spał i odpoczywał:}%

\latN{\mred{10.}~Quóniam tu, Dómine, singuláriter \textbf{in} spe~\red{*}  con\underline{\textit{sti}}\underline{\textit{tu}}\textbf{ís}ti me.}\polPs{Bo ty, Panie, osobliwie w nadziei~\red{*}  postanowiłeś mię.}%

\latN{\mred{11.}~Glória Patri, et \textbf{Fí}lio,~\red{*}  et Spirí\underline{\textit{tu}}\underline{\textit{i}} \textbf{Sanc}to.}\polPs{Chwała Ojcu i Synowi,~\red{*}  i Duchowi Świętemu.}%

\latN{\mred{12.}~Sicut erat in princípio, et nunc, et \textbf{sem}per,~\red{*}  et in sǽcula sæcu\underline{\textit{ló}}\underline{\textit{rum}}. \textbf{A}men.}\polPs{Jak była na początku, teraz i zawsze,~\red{*}  i na wieki wieków. Amen.}%
\EndParallel
\subsekcja{Psalm 90}
\vspace{-2mm}
\gresetinitiallines{0}\gabcsnippet{(c4)<c><v>\footnotesize{1. }</v></c>Qui(g) há(h)bi(j)tat(j) in(j) ad(j)ju(j)tó(j)ri(j)o(j) Al(j)<b>tís</b>(k)si(jr)mi,(j.) <c><v>*</v></c> (:) in(j) pro(j)tec(j)ti(j)ó(j)ne(j) De(j)i(j) cæ(j)li(j) <i><ul>com</i></ul>(i)<i><ul>mo</i></ul>(j)<b>rá</b>(h)bi(gr)tur.(g.) (::)
<i>Flex:</i>()  in(j) di(j)e, <c><v>†</v></c>(h. j j j  ::)}
\vspace{-5mm}
\tlumaczR{Kto mieszka w spomożeniu Najwyższego,~\red{*}  w obronie Boga niebieskiego będzie przebywał.}\segmentspace
\columnratio{0.7}
\setlength{\columnseprule}{0.4pt}
\setlength{\columnsep}{.02\textwidth}
\BeginParallel
\latN{\mred{2.}~Dicet Dómino: Suscéptor meus es tu, et refúgium \textbf{me}um:~\red{*}  Deus meus sperá\underline{\textit{bo}} \underline{\textit{in}} \textbf{e}um.}\polPs{Rzecze Panu: Obrońca mój, i ucieczka moja jesteś ty:~\red{*}  Bóg mój, w nim będę miał nadzieję.}%

\latN{\mred{3.}~Quóniam ipse liberávit me de láqueo ve\textbf{nán}tium,~\red{*}  et a \underline{\textit{ver}}\underline{\textit{bo}} \textbf{ás}pero.}\polPs{Albowiem on mię wyrwał z sidła łowiących,~\red{*}  i od przykrego słowa.}%

\latN{\mred{4.}~Scápulis suis obumbrábit \textbf{ti}bi:~\red{*}  et sub pennis e\underline{\textit{jus}} \underline{\textit{spe}}\textbf{rá}bis.}\polPs{Plecami swemi okryje cię:~\red{*}  a pod skrzydłami jego nadzieję mieć będziesz.}%

\latN{\mred{5.}~Scuto circúmdabit te véritas \textbf{e}jus:~\red{*}  non timébis a timó\underline{\textit{re}} \underline{\textit{noc}}\textbf{túr}no.}\polPs{Tarczą ogarnie cię prawda jego:~\red{*}  nie ulękniesz się od strachu nocnego.}%

\latN{\mred{6.}~A sagítta volánte in die,~\red{†} a negótio perambulánte in \textbf{té}nebris:~\red{*}  ab incúrsu et dæmónio me\underline{\textit{ri}}\underline{\textit{di}}\textbf{á}no.}\polPs{Od strzały lecącéj we dnie,~\red{†} od sprawy chodzącéj w ciemności:~\red{*}  od najazdu i od czarta południowego.}%

\latN{\mred{7.}~Cadent a látere tuo mille,~\red{†} et decem míllia a dextris \textbf{tu}is:~\red{*}  ad te autem non ap\underline{\textit{pro}}\underline{\textit{pin}}\textbf{quá}bit.}\polPs{Padnie po boku twoim tysiąc,~\red{†} a dziesięć tysięcy po prawéj stronie twojéj:~\red{*}  a ku tobie się nie przybliży.}%

\latN{\mred{8.}~Verúmtamen óculis tuis conside\textbf{rá}bis:~\red{*}  et retributiónem peccató\underline{\textit{rum}} \underline{\textit{vi}}\textbf{dé}bis.}\polPs{Zaiste oczyma twemi oglądasz:~\red{*}  a ujrzysz odpłatę grzeszników.}%

\latN{\mred{9.}~Quóniam tu es, Dómine, spes \textbf{me}a:~\red{*}  Altíssimum posuísti refú\underline{\textit{gi}}\underline{\textit{um}} \textbf{tu}um.}\polPs{Albowiem ty, Panie, jesteś nadzieja moja:~\red{*}  postawiłeś Najwyższego ucieczką twoją.}%

\latN{\mred{10.}~Non accédet ad te \textbf{ma}lum:~\red{*}  et flagéllum non appropinquábit taberná\underline{\textit{cu}}\underline{\textit{lo}} \textbf{tu}o.}\polPs{Nie przystąpi do ciebie złe:~\red{*}  i nie przybliży się bicz do przybytku twego.}%

\latN{\mred{11.}~Quóniam Angelis suis mandávit \textbf{de} te:~\red{*}  ut custódiant te in ómnibus \underline{\textit{vi}}\underline{\textit{is}} \textbf{tu}is.}\polPs{Albowiem Aniołom swoim rozkazał o tobie:~\red{*}  aby cię strzegli na wszystkich drogach twoich.}%

\latN{\mred{12.}~In mánibus por\textbf{tá}bunt te:~\red{*}  ne forte offéndas ad lápidem \underline{\textit{pe}}\underline{\textit{dem}} \textbf{tu}um.}\polPs{Na ręku będą cię nosić:~\red{*}  byś snadź nie obraził o kamień nogi twojéj.}%

\latN{\mred{13.}~Super áspidem et basilíscum ambu\textbf{lá}bis:~\red{*}  et conculcábis leónem \underline{\textit{et}} \underline{\textit{dra}}\textbf{có}nem.}\polPs{Po źmii i po bazyliszku chodzić będziesz:~\red{*}  i podepczesz lwa i smoka.}%

\latN{\mred{14.}~Quóniam in me sperávit, liberábo \textbf{e}um:~\red{*}  prótegam eum quóniam cognóvit \underline{\textit{no}}\underline{\textit{men}} \textbf{me}um.}\polPs{Iż we mnie nadzieję miał, wybawię go:~\red{*}  obronię go, że poznał imię moje.}%

\latN{\mred{15.}~Clamábit ad me, et ego exáudiam eum:~\red{†} cum ipso sum in tribulati\textbf{ó}ne:~\red{*}  erípiam eum et glorifi\underline{\textit{cá}}\underline{\textit{bo}} \textbf{e}um.}\polPs{Będzie wołał ku mnie, a ja go wysłucham:~\red{†} z nim jestem w ucisku:~\red{*} wyrwę go i uwielbię go.}%

\latN{\mred{16.}~Longitúdine diérum replébo \textbf{e}um:~\red{*}  et osténdam illi salu\underline{\textit{tá}}\underline{\textit{re}} \textbf{me}um.}\polPs{Długością dni napełnię go:~\red{*}  a okażę mu zbawienie moje.}%

\latN{\mred{17.}~Glória Patri, et \textbf{Fí}lio,~\red{*}  et Spirí\underline{\textit{tu}}\underline{\textit{i}} \textbf{Sanc}to.}\polPs{Chwała Ojcu i Synowi,~\red{*}  i Duchowi Świętemu.}%

\latN{\mred{18.}~Sicut erat in princípio, et nunc, et \textbf{sem}per,~\red{*}  et in sǽcula sæcu\underline{\textit{ló}}\underline{\textit{rum}}. \textbf{A}men.}\polPs{Jak była na początku, teraz i zawsze,~\red{*}  i na wieki wieków. Amen.}%
\EndParallel
\subsekcja{Psalm 133}
\vspace{-2mm}
\gresetinitiallines{0}\gabcsnippet{(c4)<c><v>\footnotesize{1. }</v></c>Ec(g)ce(h) nunc(j) be(j)ne(j)dí(j)ci(j)te(j) <b>Dó</b>(k)mi(jr)num,(j.) <c><v>*</v></c> (:) om(j)nes(j) <i><ul>ser</i></ul>(i)<i><ul>vi</i></ul>(j) <b>Dó</b>(h)mi(gr)ni:(g.) (::)}
\tlumaczR{Otóż teraz błogosławcie Pana,~\red{*}  wszyscy słudzy Pańscy:}\segmentspace
\columnratio{0.7}
\setlength{\columnseprule}{0.4pt}
\setlength{\columnsep}{.02\textwidth}
\BeginParallel
\latN{\mred{2.}~Qui statis in domo \textbf{Dó}mini,~\red{*}  in átriis domus \underline{\textit{De}}\underline{\textit{i}} \textbf{nos}tri.}\polPs{Którzy stoicie w domu Pańskim,~\red{*}  w sieniach domu Boga naszego.}%

\latN{\mred{3.}~In nóctibus extóllite manus vestras in \textbf{sanc}ta,~\red{*}  et benedí\underline{\textit{ci}}\underline{\textit{te}} \textbf{Dó}minum.}\polPs{W nocy podnoście ręce wasze ku świątnicy,~\red{*}  a błogosławcie Pana.}%

\latN{\mred{4.}~Benedícat te Dóminus ex \textbf{Si}on,~\red{*}  qui fecit cæ\underline{\textit{lum}} \underline{\textit{et}} \textbf{ter}ram.}\polPs{Niechaj cię błogosławi Pan z Syonu,~\red{*}  który stworzył niebo i ziemię.}%

\latN{\mred{5.}~Glória Patri, et \textbf{Fí}lio,~\red{*}  et Spirí\underline{\textit{tu}}\underline{\textit{i}} \textbf{Sanc}to.}\polPs{Chwała Ojcu i Synowi,~\red{*}  i Duchowi Świętemu.}%

\latN{\mred{6.}~Sicut erat in princípio, et nunc, et \textbf{sem}per,~\red{*}  et in sǽcula sæcu\underline{\textit{ló}}\underline{\textit{rum}}. \textbf{A}men.}\polPs{Jak była na początku, teraz i zawsze,~\red{*}  i na wieki wieków. Amen.}%
\EndParallel
\newpage
 \small\emph{Per Annum}\linespace
\gresetinitiallines{1}\greannotation{Ant.}\greannotation{8. G}\gabcsnippet{(c4) MI(g)se(g')ré(h)re(gf) *() mi(gh)hi(h') Dó(g)mi(fg)ne,(g.) (;) et(h) ex(j)áu(i_[uh:l]j)di(h) o(h)ra(hg)ti(f)ó(gh)nem(h) me(g.)am.(g.) (::)}
\tlumaczR{Zmiłuj się nademną, Panie, a wysłuchaj modlitwę moją.}
\small\emph{Tempore Paschali}\linespace
\gresetinitiallines{1}\greannotation{Ant.}\greannotation{8. G}\gabcsnippet{(c4) AL(g)le-(g')lú(h)ia,(gf__) (,) al(h)le-(j)lú-(i_j)ia,(h_GF) al(gh)le-(h)lú-(g.)ia.(g.) (::)}
\sekcja{Modlitwy końcowe}{Odmawia się po cichu w pozycji klęczącej.}
\sloppy
\wstawplik{modlitwy_koncowe_credo}

\end{document}